% !TEX root = ../main.tex
% ---------------------------------------------------------------
% GENERATED FILE - DO NOT EDIT.
% Produced by docs/tools/render_reference.py from the repository source.
% Regenerate with: make docs
% ---------------------------------------------------------------

\apibreak
\section{\texttt{ExperimentConfig}}
\label{class:feedback_intelligence_agent.experiments.ExperimentConfig}\label{schema:ExperimentConfig}\index{Classes!ExperimentConfig}

Typed description of one experiment run.
\apifield{Inherits}\texttt{BaseModel}
\apifield{Fields}
\begin{arglist}
  \item[\texttt{name}] \texttt{str}, default \texttt{'experiment'}, \texttt{min\_\allowbreak{}length=1}
  \item[\texttt{description}] \texttt{str}, default \texttt{''}
  \item[\texttt{dataset\_\allowbreak{}path}] \texttt{Path}, default \texttt{Path('data/\allowbreak{}sample\_\allowbreak{}feedback.\allowbreak{}csv')\allowbreak{}}
  \item[\texttt{queries\_\allowbreak{}path}] \texttt{Path}, default \texttt{Path('examples/\allowbreak{}queries.\allowbreak{}jsonl')\allowbreak{}}
  \item[\texttt{output\_\allowbreak{}dir}] \texttt{Path}, default \texttt{Path('.\allowbreak{}artifacts/\allowbreak{}experiments/\allowbreak{}default')\allowbreak{}}
  \item[\texttt{chunk\_\allowbreak{}size}] \texttt{int}, default \texttt{80}, \texttt{ge=1}. Maximum words per chunk.
  \item[\texttt{chunk\_\allowbreak{}overlap}] \texttt{int}, default \texttt{16}, \texttt{ge=0}. Words shared between chunks.
  \item[\texttt{top\_\allowbreak{}k}] \texttt{int}, default \texttt{4}, \texttt{ge=1,\allowbreak{} le=12}
  \item[\texttt{embedding\_\allowbreak{}provider}] \texttt{Literal['hashing']\allowbreak{}}, default \texttt{'hashing'}
  \item[\texttt{embedding\_\allowbreak{}dim}] \texttt{int}, default \texttt{512}, \texttt{ge=64,\allowbreak{} le=8192}
  \item[\texttt{llm\_\allowbreak{}provider}] \texttt{Literal['local',\allowbreak{} 'openai',\allowbreak{} 'openai\_\allowbreak{}responses',\allowbreak{} 'bedrock']\allowbreak{}}, default \texttt{'local'}
  \item[\texttt{retriever\_\allowbreak{}type}] \texttt{RetrieverType}, default \texttt{'dense'}
  \item[\texttt{dense\_\allowbreak{}weight}] \texttt{float}, default \texttt{0.\allowbreak{}6}, \texttt{ge=0.\allowbreak{}0}
  \item[\texttt{lexical\_\allowbreak{}weight}] \texttt{float}, default \texttt{0.\allowbreak{}4}, \texttt{ge=0.\allowbreak{}0}
  \item[\texttt{tracking\_\allowbreak{}provider}] \texttt{TrackingProvider}, default \texttt{'none'}
  \item[\texttt{tracking\_\allowbreak{}uri}] \texttt{str \textbar{} None}, default \texttt{None}
  \item[\texttt{tracking\_\allowbreak{}experiment\_\allowbreak{}name}] \texttt{str}, default \texttt{'feedback-\allowbreak{}intelligence-\allowbreak{}agent'}
\end{arglist}
\apifield{Details}The configuration covers the full pipeline: chunking, embedding, retrieval strategy, answer generation, and the dataset and query files used for evaluation.
\apifield{Source}\texttt{src/\allowbreak{}feedback\_\allowbreak{}intelligence\_\allowbreak{}agent/\allowbreak{}experiments.\allowbreak{}py:52}~(\href{https://github.com/DiogoRibeiro7/feedback-intelligence-agent/blob/936886da072fa7025a57c5517b9c90772bedae09/src/feedback_intelligence_agent/experiments.py#L52}{online}), pydantic model, module \cref{module:feedback_intelligence_agent.experiments}

\subsection{\texttt{check\_chunk\_overlap}}
\label{method:feedback_intelligence_agent.experiments.ExperimentConfig.check_chunk_overlap}\index{Methods!check\_chunk\_overlap}
Reject overlaps that would prevent chunking from terminating.
\apifield{Usage}
\begin{lstlisting}[style=usage, language=Python]
check_chunk_overlap(self) -> ExperimentConfig
\end{lstlisting}
\apifield{Value}\texttt{ExperimentConfig}
\apifield{Raises}\texttt{ValueError}
\apifield{Source}\texttt{src/\allowbreak{}feedback\_\allowbreak{}intelligence\_\allowbreak{}agent/\allowbreak{}experiments.\allowbreak{}py:79}~(\href{https://github.com/DiogoRibeiro7/feedback-intelligence-agent/blob/936886da072fa7025a57c5517b9c90772bedae09/src/feedback_intelligence_agent/experiments.py#L79}{online})

\subsection{\texttt{from\_yaml}}
\label{method:feedback_intelligence_agent.experiments.ExperimentConfig.from_yaml}\index{Methods!from\_yaml}
Load and validate an experiment configuration from a YAML file.
\apifield{Usage}
\begin{lstlisting}[style=usage, language=Python]
from_yaml(cls, path: str | Path) -> ExperimentConfig
\end{lstlisting}
\apifield{Arguments}
\begin{arglist}
  \item[\texttt{path}] \texttt{str \textbar{} Path}, required.
\end{arglist}
\apifield{Value}\texttt{ExperimentConfig}
\apifield{Raises}\texttt{ValueError}
\apifield{Source}\texttt{src/\allowbreak{}feedback\_\allowbreak{}intelligence\_\allowbreak{}agent/\allowbreak{}experiments.\allowbreak{}py:86}~(\href{https://github.com/DiogoRibeiro7/feedback-intelligence-agent/blob/936886da072fa7025a57c5517b9c90772bedae09/src/feedback_intelligence_agent/experiments.py#L86}{online})
